% !TeX root = ../../Book1.tex
% 纲要标题：## 第9章　有限生成阿贝尔群
% 纲要位置：计划纲要.md，第 356 行。

\chapter{有限生成阿贝尔群}
\label{b1:ch:09}
\BookChapterNumber{b1:ch:09}

\begin{chapterabstract}
有限生成阿贝尔群可以完全分类。本章从整数格出发，证明自由阿贝尔群的子群仍自由，分离挠与无挠部分，再通过整数关系矩阵的 Smith 化简建立不变因子分解。唯一性由各素数的挠层数恢复，不借助后续模论。最后给出商格、指数和坐标变换的具体计算。
\end{chapterabstract}
\begin{learninggoals}
\begin{itemize}
\item 分清整数线性无关、自由基、有限生成与无挠之间的关系。
\item 会把有限阿贝尔群在初等因子与不变因子形式之间转换，并解释唯一性。
\item 能由整数关系矩阵计算商群，保留需要追踪元素时的坐标变换。
\end{itemize}
\end{learninggoals}
本章仅需前面群同态、商群、循环群的基础，以及整数的带余除法、Bézout 等式和线性代数中的行列式。阿贝尔群统一使用加法。一般无限秩子群自由的证明使用良序与选择公理；以有限生成分类为目标时，只需该证明的有限秩部分。有关生成关系与结构定理的伴读可参阅 Dummit 与 Foote 的《Abstract Algebra》\cite{DummitFoote2003}，本章的初等证明不依赖第二卷的模论。

% !TeX root = ../../Book1.tex
% 纲要标题：### 9.1 自由阿贝尔群
% 纲要位置：计划纲要.md，第 358 行。

\section{自由阿贝尔群}
\label{b1:sec:0901}

% 纲要内容（仅作写作计划，尚非教材正文）：
%
% * 整数格与自由阿贝尔群
% * 子群仍自由及有限秩情形
% * 秩及其不变性
%

\subsection{整数格与自由阿贝尔群}
\label{b1:subsec:0901:01}

阿贝尔群的乘法可以改写为加法，生成元之间的关系因而成为整数线性关系。本章统一用 \(0\) 表示单位元，用 \(mx\) 表示元素 \(x\) 的整数倍。先从没有非平凡整数关系的群开始。

\begin{definition}\label{b1:def:free-abelian}
设 \(I\) 为集合，令
\[
\mathbb Z^{(I)}
=\{\,(n_i)_{i\in I}\mid n_i\in\mathbb Z,\ %
n_i\ne0\text{ 的指标只有有限个}\,\}.
\]
逐分量加法使它成为阿贝尔群。与某个 \(\mathbb Z^{(I)}\) 同构的群称为\term[ziyou-abeier-qun]{自由阿贝尔群}[free abelian group]。其中标准向量 \(e_i\) 的对应元素构成一组基：每个元素均可唯一写为有限和 \(\sum n_i e_i\)。
\end{definition}
\symidx{\(\mathbb Z^{(I)}\)：有限支集整数族}
有限的 \(I\) 给出熟悉的整数格 \(\mathbb Z^n\)。无限指标时必须限定有限支集；允许任意整数族所得的全直积是不同的构造，本章不把它当作自由阿贝尔群的定义。

\begin{proposition}\label{b1:prop:free-abelian-universal}
对任意阿贝尔群 \(A\) 和映射 \(f:I\to A\)，存在唯一同态
\(\widetilde f:\mathbb Z^{(I)}\to A\)，使 \(\widetilde f(e_i)=f(i)\)。
\end{proposition}
\begin{proof}
定义 \(\widetilde f(\sum n_i e_i)=\sum n_i f(i)\)。表示的唯一性保证良定义；和只有有限项，交换性保证加法相容。生成元的像已经指定，所以这样的同态唯一。
\end{proof}
这与自由群的泛性质相似，区别在于目标限为阿贝尔群。比如从 \(\mathbb Z^2\) 到群 \(G\) 指定两个基向量的像 \(x,y\)，必须要求 \(xy=yx\)；从 \(F_2\) 出发则没有这个要求。

\ProseParagraphBreak
在 \(\mathbb Z^2\) 中，\((1,0),(1,1)\) 构成一组基，因为
\((u,v)=(u-v)(1,0)+v(1,1)\)，系数唯一。而 \((2,0),(0,1)\) 虽然整数线性无关，却只生成第一坐标为偶数的子群。整数格的“无关”与“生成”必须分别检查，不能把实向量空间中的一组基自动当作整数格的基。

\subsection{自由阿贝尔群的子群}
\label{b1:subsec:0901:02}

\begin{theorem}\label{b1:thm:free-abelian-subgroup}
自由阿贝尔群的每个子群仍是自由阿贝尔群。若 \(H\le\mathbb Z^n\)，则 \(H\) 有一组至多 \(n\) 个元素的基。
\end{theorem}
\begin{proof}
先证明有限秩情形，对 \(n\) 归纳。\(n=0\) 显然。对 \(n\ge1\)，把 \(H\) 投影到最后一个坐标，其像是 \(\mathbb Z\) 的子群，所以为 \(d\mathbb Z\)，其中 \(d\ge0\)。核
\[
K=H\cap(\mathbb Z^{n-1}\times\{0\})
\]
由归纳假设自由且秩至多 \(n-1\)。若 \(d=0\)，则 \(H=K\)。若 \(d>0\)，选 \(h\in H\) 的末坐标为 \(d\)。任意 \(x\in H\) 的末坐标为 \(md\)，于是 \(x-mh\in K\)；而 \(mh\in K\) 迫使 \(md=0\)，进而 \(m=0\)。所以 \(H=K\oplus\mathbb Zh\)，在 \(K\) 的基后加上 \(h\) 即可。

一般情形在通常的选择公理下，把母群的一组基良序为 \((e_\alpha)_{\alpha<\kappa}\)。令 \(F_\alpha\) 由所有 \(e_\beta\)（\(\beta<\alpha\)）生成，令 \(H_\alpha=H\cap F_\alpha\)。从 \(H_{\alpha+1}\) 取 \(e_\alpha\) 的坐标，像为 \(d_\alpha\mathbb Z\)。若 \(d_\alpha>0\)，选 \(h_\alpha\) 的此坐标为 \(d_\alpha\)，完全相同的减去整数倍步骤给出
\[
H_{\alpha+1}=H_\alpha\oplus\mathbb Zh_\alpha.
\]
若 \(d_\alpha=0\)，两项相等。对极限指标 \(\lambda\)，有限支集保证
\(H_\lambda=\bigcup_{\alpha<\lambda}H_\alpha\)。

所选全部 \(h_\alpha\) 生成 \(H\)：逐级使用上式，并在极限步骤取并即可；任何给定元素只涉及有限项。它们也整数线性无关，因为任何有限关系中指标最大的 \(h_\alpha\) 在 \(e_\alpha\) 位置的系数为其整数系数乘 \(d_\alpha\)，更早的项在该位置全为零，所以该系数必须为零；依次向前得到所有系数为零。这便是一组自由基。
\end{proof}
证明并未断言子群的基总能扩充为母群的基。子群 \(2\mathbb Z\le\mathbb Z\) 的基为 \(2\)，却不能成为 \(\mathbb Z\) 的基的一部分；商群在这里出现了二阶挠元。作为可跟算的二维例子，
\[
H=\{\,(u,v)\in\mathbb Z^2\mid u+v\equiv0\pmod3\,\}
\]
在第二坐标的投影为 \(\mathbb Z\)，可选 \(h=(-1,1)\)；核为 \(3\mathbb Z\times\{0\}\)。所以 \((3,0),(-1,1)\) 是 \(H\) 的一组基。

\subsection{秩及其不变性}
\label{b1:subsec:0901:03}

基的个数必须与选择无关，才可用来比较群。

\begin{proposition}\label{b1:prop:free-abelian-rank}
自由阿贝尔群任意两组基的基数相同，称为\term[ziyou-abeier-qun-de-zhi]{自由阿贝尔群的秩}[rank of a free abelian group]。
\end{proposition}
\begin{proof}
若一组基有有限个 \(n\) 元，则每个元素模 \(2A=\{2x\mid x\in A\}\) 后，可唯一把各坐标减为 \(0\) 或 \(1\)，所以 \(|A/2A|=2^n\)。另一组有限基因而具有相同大小。若 \(A\) 有有限生成集，则这些生成元相对于任意一组基只涉及有限多个基元素；其余基元素不可能由它们生成，所以每组基都有限。

对两组无限基 \(B,C\)，把 \(C\) 中的每个元素写为 \(B\) 的有限线性组合。这些支集的并必须覆盖 \(B\)，故 \(|B|\le |C|\)；交换角色得到反向不等式。这与自由群的无限秩证明使用相同的有限支集和基数比较。
\end{proof}
整数线性无关最多只能给出秩的下界。例如前面的指数三子群与 \(\mathbb Z^2\) 都有秩二，但并不相等。另一方面，秩为一的自由阿贝尔群必为无限循环群，因此加法群 \(\mathbb Q\) 若自由就不可能仅凭“嵌在一条实直线中”断言它秩为一；下一节会直接证明它根本不自由。

\ProseParagraphBreak
有限生成阿贝尔群都有从某个 \(\mathbb Z^n\) 出发的满同态。其核由本节定理仍是有限秩自由阿贝尔群，于是关系可以选为有限多个整数向量。这是后面使用有限整数矩阵进行分类的依据；有限生成的关系并非一开始就应当默认为已知。

% !TeX root = ../../Book1.tex
% 纲要标题：### 9.2 挠与无挠
% 纲要位置：计划纲要.md，第 364 行。

\section{挠与无挠}
\label{b1:sec:0902}

% 纲要内容（仅作写作计划，尚非教材正文）：
%
% * 挠子群与商群
% * 有限阿贝尔群的 p 主分解
% * 有限生成条件不可省略的例子
%

\subsection{挠子群与商群}
\label{b1:subsec:0902:01}

整数格中非零向量没有有限阶，但一般阿贝尔群可以同时含有无限阶和有限阶元素。先把后一部分提取出来，可以把后续分类拆成自由方向与有限阶方向。

\begin{definition}\label{b1:def:torsion-subgroup}
阿贝尔群 \(A\) 中满足某个正整数 \(m\) 使 \(mx=0\) 的元素称为\term[naoyuan]{挠元}[torsion element]。全部挠元组成的集合记为
\[
A_{\mathrm{tor}}\coloneqq\{\,x\in A\mid mx=0
\text{ 对某个整数 }m\ge1\,\}.
\]
若 \(A_{\mathrm{tor}}=A\)，称为\term[naoqun]{挠群}[torsion group]；若 \(A_{\mathrm{tor}}=0\)，称为\term[wunao-qun]{无挠群}[torsion-free group]。
\end{definition}
\symidx{\(A_{\mathrm{tor}}\)：挠子群}

\begin{proposition}\label{b1:prop:torsion-quotient}
\(A_{\mathrm{tor}}\) 是 \(A\) 的特征子群，称为\term[naoziqun]{挠子群}[torsion subgroup]；商群 \(A/A_{\mathrm{tor}}\) 无挠。
\end{proposition}
\begin{proof}
若 \(mx=0,ny=0\)，则交换性给出 \(mn(x-y)=0\)，故挠元对差封闭。同态保持整数倍，所以自同构保持挠元，给出特征性。若 \(k(x+A_{\mathrm{tor}})=0\)，则 \(kx\) 为挠元，存在 \(m\ge1\) 使 \(mkx=0\)。于是 \(x\in A_{\mathrm{tor}}\)，原陪集就是零陪集。
\end{proof}
交换性不可随意去掉。在无限二面体群中，两次反射都有二阶，而其乘积可以是无限阶平移。因此一般群的有限阶元素集合未必为子群。

\begin{proposition}\label{b1:prop:fg-abelian-subgroups}
有限生成阿贝尔群的每个子群都有限生成；其中挠子群为有限群。
\end{proposition}
\begin{proof}
取满同态 \(\pi:\mathbb Z^n\to A\)。对子群 \(B\le A\)，原像 \(\pi^{-1}(B)\le\mathbb Z^n\) 由前节定理有有限基，其像生成 \(B\)。特别地，\(A_{\mathrm{tor}}\) 有有限生成元 \(t_1,\ldots,t_s\)。若 \(t_i\) 的阶为 \(m_i\)，每个挠元都能写成
\(\sum a_i t_i\)，其中 \(0\le a_i<m_i\)，故元素数至多 \(\prod m_i\)，为有限群。
\end{proof}
在 \(A=\mathbb Z\oplus C_6\) 中，\((u,\bar v)\) 为挠元当且仅当 \(u=0\)，所以挠子群为 \(C_6\)，无挠商群为 \(\mathbb Z\)。一般有限生成阿贝尔群是否都能像这样拆成直和，仍需下一节的结构定理。

\subsection{有限阿贝尔群的主分解}
\label{b1:subsec:0902:02}

不同素数的挠可以独立拆开。对素数 \(p\)，定义
\[
A_p=\{\,x\in A\mid p^k x=0\text{ 对某个 }k\ge0\,\},
\]
称为 \(A\) 的\term[p-zhufenliang]{\(p\) 主分量}[p-primary component]。它是特征子群，因为两个元素可取同一个较大的 \(p\) 幂来消去，且同态保持其消去关系。
\symidx{\(A_p\)：阿贝尔群的 \(p\) 主分量}

\begin{theorem}\label{b1:thm:abelian-primary}
设 \(A\) 为有限阿贝尔群，\(|A|=p_1^{e_1}\cdots p_t^{e_t}\)。则加法映射给出同构
\[
A_{p_1}\oplus\cdots\oplus A_{p_t}\ \cong\ A.
\]
\end{theorem}
\begin{proof}
若 \(A=0\)，素因子列表为空，空直和为零群，结论成立。以下设 \(A\ne0\)。令 \(N=|A|\)、\(N_i=N/p_i^{e_i}\)。这些 \(N_i\) 的最大公因数为一，故 Bézout 等式给出整数 \(c_i\) 使 \(\sum_i c_iN_i=1\)。对任意 \(x\in A\)，Lagrange 定理给出 \(Nx=0\)，所以
\[
x=\sum_i x_i,\qquad x_i=c_iN_i x,\qquad
p_i^{e_i}x_i=0.
\]
这说明各主分量之和为 \(A\)。

为证唯一性，设 \(\sum_i y_i=0\)，其中 \(y_i\in A_{p_i}\)。每个 \(y_i\) 的阶既为 \(p_i\) 幂又整除 \(N\)，所以 \(p_i^{e_i}y_i=0\)。对等式施以整数倍 \(c_jN_j\)：当 \(i\ne j\) 时，\(N_j\) 含因子 \(p_i^{e_i}\)，故对应项为零；当 \(i=j\) 时，由 Bézout 等式模 \(p_j^{e_j}\) 得 \(c_jN_j\equiv1\)，对应项仍为 \(y_j\)。于是 \(y_j=0\)，对每个 \(j\) 都成立。
\end{proof}
这称为有限阿贝尔群的\term[zhufenjie]{主分解}[primary decomposition]。例如在 \(C_{12}\) 中，对 \(N_2=3,N_3=4\)，可取 \(c_2=-1,c_3=1\)。所以每个 \(\bar x\) 唯一写为
\[
\bar x=(-3\bar x)+(4\bar x),
\]
前项被 \(4\) 消去，后项被 \(3\) 消去，对应 \(C_{12}\cong C_4\oplus C_3\)。这种投影是由群上的整数倍映射给出的，不需要预先选择每个主分量的生成元。

\subsection{有限生成条件与反例}
\label{b1:subsec:0902:03}

有限生成条件的作用至少有两处：它让关系由有限矩阵描述，也让挠子群有限。两个具体例子可以防止把有限生成结构定理误用于一般阿贝尔群。

\ProseParagraphBreak
首先，\((\mathbb Q,+)\) 无挠，但不有限生成。若有限多个有理数的分母都整除 \(N\)，它们的整数线性组合都在 \(\frac1N\mathbb Z\) 中；有理数 \(\frac1{2N}\) 不在其中。它也不是自由阿贝尔群：在 \(\mathbb Q\) 中每个元素都可除以二，故 \(\mathbb Q=2\mathbb Q\)；在任何非零自由阿贝尔群中，一个基元素都不属于二倍子群。这个区别排除了任意有限或无限秩的自由基。

\ProseParagraphBreak
其次，固定素数 \(p\)，考虑
\[
C_{p^\infty}
=\mathbb Z[1/p]/\mathbb Z,\qquad
\mathbb Z[1/p]=\{\,a/p^k\mid a\in\mathbb Z,\ k\ge0\,\}.
\]
这里 \(\mathbb Z[1/p]\) 只作为有理数加法的子群使用。商群中 \(\frac1{p^k}+\mathbb Z\) 的阶恰为 \(p^k\)，所以它是无限挠群。有限多个元素都在某一个由 \(\frac1{p^k}+\mathbb Z\) 生成的有限循环群中，因而整个群不有限生成。这个群称为\term[prufer-qun]{Prüfer 群}[Prüfer group]。

\ProseParagraphBreak
上述两个群还共享一种现象：有限生成的小部分很简单，全部小部分的并却未必具有同样的有限描述。因此后续每次使用“有限个不变因子”或“挠子群有限”时，都应保留有限生成这一条件。下一节证明的无挠推出自由，也明确限于有限生成阿贝尔群。

% !TeX root = ../../Book1.tex
% 纲要标题：### 9.3 结构定理
% 纲要位置：计划纲要.md，第 370 行。

\section{结构定理}
\label{b1:sec:0903}

% 纲要内容（仅作写作计划，尚非教材正文）：
%
% * 不变因子分解与初等因子分解
% * 结构定理的初等证明
% * 分解的唯一性
%

\subsection{不变因子分解与初等因子分解}
\label{b1:subsec:0903:01}

本节的目标是把任意有限生成阿贝尔群化为循环群的有限直和，而且给出可由群本身恢复的分类数据。有限循环因子有两种常用排列方式：按整除关系排列，或按素数幂拆开。

\begin{theorem}\label{b1:thm:fga-structure}
每个有限生成阿贝尔群 \(A\) 都有分解
\[
A\cong\mathbb Z^r\oplus C_{d_1}\oplus\cdots\oplus C_{d_s},
\qquad 1<d_1\mid d_2\mid\cdots\mid d_s,
\]
其中 \(r,s\ge0\)。整数 \(r\) 及序列 \(d_1,\ldots,d_s\) 由 \(A\) 唯一确定。也可以唯一地写为一个自由部分与若干素数幂阶循环群的直和，唯一性允许重新排列这些循环因子。
\end{theorem}
第一种称为\term[bubian-yinzi-fenjie]{不变因子分解}[invariant factor decomposition]，其中 \(d_i\) 为不变因子；第二种称为\term[chudeng-yinzi-fenjie]{初等因子分解}[elementary divisor decomposition]，各循环因子的素数幂阶称为初等因子。下文依次证明存在性与唯一性。

\ProseParagraphBreak
例如 \(C_4\oplus C_6\) 还没有按整除关系排列，因为 \(4\nmid6\)。先按素数拆开，
\[
C_4\oplus C_6\cong C_4\oplus C_2\oplus C_3
\cong C_2\oplus C_{12}.
\]
初等因子是 \(4,2,3\)，不变因子是 \(2,12\)。这两种写法表达同一个群，但各自必须遵守相应的归一化条件；不能在任意循环因子拆分之间直接谈唯一性。

\subsection{结构定理的初等证明}
\label{b1:subsec:0903:02}

生成元与关系构成一个整数矩阵。要整理关系，只能用可逆的整数行、列操作，不能像实数线性代数那样任意除以主元；除法会改变整数格。

\begin{lemma}[整数矩阵的 Smith 化简]\label{b1:lem:smith-integer}
对任意 \(m\times n\) 整数矩阵 \(M\)，可以通过交换两行或两列、把一行或一列乘 \(-1\)、把某行或列的整数倍加到另一行或列，化为矩形对角阵
\[
\operatorname{diag}(d_1,\ldots,d_t,0,\ldots),
\qquad 0<d_1\mid d_2\mid\cdots\mid d_t.
\]
\end{lemma}
\begin{proof}
若矩阵全零则结束。否则把一个非零元素移到左上角，并使其为正，记为 \(a\)。对第一列的某个元素 \(b\)，作整数除法 \(b=qa+r\)，其中 \(0\le r<a\)，再把第一行的 \(q\) 倍从该行减去。如果 \(r>0\)，交换该行与第一行，使正主元严格变小；若 \(r=0\)，就消去了该元素。对第一行作同样的列操作。每当出现非零余数，主元严格下降，因此这种下降只可能发生有限次；主元不下降时，有限步内可把第一行和第一列的其余元素都清零。

此时矩阵形如
\[
\begin{pmatrix}a&0\\0&B\end{pmatrix}.
\]
如果 \(a\) 不整除 \(B\) 中的某个元素 \(b_{ij}\)，把该元素所在的行加到第一行。左上角仍为 \(a\)，而第一行现在含有 \(b_{ij}\)。再次用列的整数除法就产生小于 \(a\) 的非零余数，使主元严格下降，然后重新清理首行首列。正整数不能无限下降，所以最终获得这样的分块形式，且 \(a\) 整除 \(B\) 的全部元素。

固定首行首列，对 \(B\) 递归进行相同步骤。行列数都减少，递归必停。由于 \(B\) 的所有元素一开始均为 \(a\) 的倍数，后续整数行列操作保持此性质，所以它所得的非零对角元也都是 \(a\) 的倍数。再用递归中的整除链，便得到全部 \(d_i\mid d_{i+1}\)。所用每一步操作都有整数逆操作，没有改变相应整数格的可逆坐标性质。
\end{proof}
上述结果称为整数矩阵的\term[smith-biaozhunxing]{Smith 标准形}[Smith normal form]。证明给出的步骤是一个终止算法，而不只是允许“选取适当基”的存在性叙述。

\begin{proof}[结构定理的存在性证明]
取 \(A\) 的有限生成元，得到满同态 \(\pi:\mathbb Z^m\to A\)。由\cref{b1:thm:free-abelian-subgroup}，核 \(K\) 有有限基 \(k_1,\ldots,k_t\)，其中 \(t\le m\)。把它们在标准基下的坐标作为列，组成 \(m\times t\) 整数矩阵 \(M\)。

整数可逆列操作只改变 \(K\) 的生成基，因而不改变其生成子群；整数可逆行操作是母群 \(\mathbb Z^m\) 的自同构，因而把商群变成同构的商群。对 \(M\) 作 Smith 化简后，核变成由
\[
d_1e_1,\ldots,d_te_t
\]
生成的子群。这里恰有 \(t\) 个非零对角元：原列在整数上无关，也就在有理数上无关，因为有理数关系乘以公共分母就成为整数关系；可逆行列操作保持这个性质。于是
\[
A\cong
\mathbb Z^m/K
\cong C_{d_1}\oplus\cdots\oplus C_{d_t}\oplus\mathbb Z^{m-t}.
\]
等于一的对角元只贡献平凡群，删去它们便得到不变因子形式。

最后把每个 \(d_i\) 作素因数分解。对互素的 \(u,v\)，映射
\(C_{uv}\to C_u\oplus C_v\)、\(\bar x\mapsto(\bar x,\bar x)\)
的核平凡，且两边元素数相同，所以为同构。反复应用便把每个有限循环因子拆成素数幂阶循环群，得到初等因子形式。
\end{proof}
证明同时得到：有限生成无挠阿贝尔群自由；一般有限生成阿贝尔群的挠子群可以分离成一个直和因子。不过具体选出的自由补群依赖基的选择，结构定理没有宣称这个补群唯一。

\subsection{分解的唯一性}
\label{b1:subsec:0903:03}

先恢复自由部分，再分别恢复每个素数的幂阶。挠子群 \(T=A_{\mathrm{tor}}\) 是由群内条件确定的，所以同构必把它映到另一个群的挠子群。上面的分解说明 \(A/T\cong\mathbb Z^r\)，于是\cref{b1:prop:free-abelian-rank}确定了 \(r\)。

\ProseParagraphBreak
对有限群 \(T\) 的 \(p\) 主分量 \(B=T_p\)，考虑全由整数倍定义的子群及商群
\[
B\supseteq pB\supseteq p^2B\supseteq\cdots,\qquad
p^{j-1}B/p^jB.
\]
若 \(B\cong\bigoplus_{i=1}^k C_{p^{e_i}}\)，每个因子对第 \(j\) 层的贡献在 \(e_i\ge j\) 时有 \(p\) 个元素，在 \(e_i<j\) 时只有一个。因此
\[
\bigl|p^{j-1}B/p^jB\bigr|=p^{c_j},
\qquad c_j=\#\{\,i\mid e_i\ge j\,\}.
\]
左边是由 \(B\) 本身确定的数，所以每个 \(c_j\) 都是同构不变量。恰好等于 \(j\) 的指数出现 \(c_j-c_{j+1}\) 次，因而全部 \(e_i\) 及其重数唯一。这证明初等因子的唯一性，没有引入模或张量积。

\ProseParagraphBreak
由初等因子恢复不变因子也没有歧义。若 \(T=0\)，有限循环因子为空，取 \(s=0\) 即可。以下设 \(T\ne0\)。对每个素数 \(p\)，把正指数按非降次序排列；令 \(s\) 为各素数指数列表长度的最大值，把较短列表的左端补零至长度 \(s\)。逐位置把相应的 \(p\) 幂相乘，便得到
\[
d_1\mid d_2\mid\cdots\mid d_s.
\]
反过来，任何满足整除链且 \(d_1>1\) 的不变因子分解，其每个素数的指数必非降，去掉左端的零以后恰为该素数的初等因子指数。至少有一个素数整除 \(d_1\)，所以某个列表的长度恰为 \(s\)，不会额外留下全为零的位置。由此不变因子也唯一，\cref{b1:thm:fga-structure}得证。

\ProseParagraphBreak
例如若一个二主群的各层商群阶依次为 \(2^3,2^2,2^2\)，再下一层为一，则指数至少为 \(1,2,3\) 的因子分别有 \(3,2,2\) 个。因此指数一出现一次，指数二不出现，指数三出现两次，该群为 \(C_2\oplus C_8\oplus C_8\)。只知道群的阶为 \(2^7\) 无法作出这一判断，逐层取整数倍才恢复了各个循环因子的长度。

% !TeX root = ../../Book1.tex
% 纲要标题：### 9.4 结构定理的计算与模论解释
% 纲要位置：计划纲要.md，第 383 行。

\section{结构定理的计算与模论解释}
\label{b1:sec:0904}

结构定理给出了分类形式，实际计算还须把生成元、关系和坐标变换一一对应。下面先固定关系矩阵的行列约定，再用一个完整例子说明怎样从化简结果读出循环因子并追踪具体元素。随后讨论整数格的商群、子格指数与子式判据。节末的模论解释只是对第二卷的预告，本节全部计算仍只依赖整数矩阵、子群与商群。

% 纲要内容（仅作写作计划，尚非教材正文）：
%
% * 由整数关系矩阵确定交换群
% * 整数格的商群与有限指数
% * 模论解释与后续推广预告
%
% ---
%

\subsection{整数关系矩阵与交换群}
\label{b1:subsec:0904:01}

使用结构定理计算时，第一步是固定矩阵的方向。本节一律把生成元列为坐标轴，把一条整数关系写成一个列向量。给定 \(m\times n\) 整数矩阵 \(M\)，可以用它的各列施加关系，并定义交换群
\[
A_M\coloneqq\mathbb Z^m/M\mathbb Z^n.
\]
这里 \(M\mathbb Z^n\) 是由这些列关系的全部整数线性组合组成的子群。列操作改变关系生成集，行操作改变群的生成坐标；两者都必须是可逆的整数操作。

\ProseParagraphBreak
若要用已知关系分析一个事先给定的交换群 \(A\)，还须核对关系是否完整。设 \(A\) 由 \(a_1,\ldots,a_m\) 生成，并令
\[
\pi:\mathbb Z^m\longrightarrow A,
\qquad e_i\longmapsto a_i
\]
为自然满同态。若 \(M\) 的各列只记录若干已经验证的关系，则只能得到 \(M\mathbb Z^n\subseteq\ker\pi\)，因而 \(\pi\) 诱导自然满同态
\[
\overline\pi:\mathbb Z^m/M\mathbb Z^n\longrightarrow A.
\]
只有进一步证明 \(M\mathbb Z^n=\ker\pi\)，这个满同态才是同构，Smith 化简所得结构才是 \(A\) 本身的分类。例如 \(A=C_2=\langle a\rangle\) 确实满足关系 \(4a=0\)，但只记录这一关系得到的是 \(\mathbb Z/4\mathbb Z\cong C_4\)，它只自然满射到 \(C_2\)；遗漏的关系 \(2a=0\) 正说明所列关系并不完整。

\ProseParagraphBreak
Smith 计算本身的输入是一张有限整数矩阵 \(M\)。在 \(\mathbb Z\) 上逐次作可逆行列操作，并同时累计这些操作，输出 \(m\times m\) 与 \(n\times n\) 整数可逆矩阵 \(U,V\)，以及 Smith 标准形 \(D=UMV\)。只求同构类型时，从 \(D\) 读取自由秩与非平凡不变因子即可；还要追踪原生成元或具体元素时，则须保留 \(U,V\) 及其逆矩阵。\cref{b1:lem:smith-integer}的证明给出了这一过程的终止性，但本章不讨论它的复杂度。若 \(M\) 用来分析一个事先给定的群，求得生成完整关系核的矩阵是 Smith 化简之前另须解决的问题。

\ProseParagraphBreak
设 Smith 标准形有 \(t\) 个非零对角元 \(d_1\mid\cdots\mid d_t\)。\cref{b1:tab:smith-coordinate-types}区分新坐标中的三种情况：前 \(t\) 个坐标受到非零关系约束，其余 \(m-t\) 个坐标没有这样的约束。

\begin{table}[htbp]
\centering
\caption{Smith 标准形中的坐标与商群因子}
\label{b1:tab:smith-coordinate-types}
\begin{tabularx}{\linewidth}{@{}>{\raggedright\arraybackslash}p{0.44\linewidth}>{\raggedright\arraybackslash}X@{}}
\toprule
新坐标中的情形 & 对商群的贡献 \\
\midrule
非单位主元 \(d_i>1\) & 非平凡有限循环因子 \(C_{d_i}\) \\
单位主元 \(d_i=1\) & 该坐标在商群中为零，平凡因子 \(C_1\) 删去 \\
没有非零主元约束的坐标 & 每个贡献一个无限循环因子，共有 \(m-t\) 个 \\
\bottomrule
\end{tabularx}
\end{table}

矩阵的非零 Smith 对角元允许为 \(1\)；群结构定理所列的不变因子则删去这些 \(1\)，只保留大于 \(1\) 的项。标准形中的零列表示经过可逆整数重组后出现的零关系，可以在这一阶段删去；这不意味着原关系组中必有某条关系可以直接删除。关系矩阵可以有任意多个相互依赖的列；自由部分的秩由目标格的坐标数减去非零主元数得到，为 \(m-t\)，不是 \(n-t\)。

\ProseParagraphBreak
例如对一生成元群施加关系 \(2x=0\)、\(3x=0\)，关系矩阵为 \(M=\begin{pmatrix}2&3\end{pmatrix}\)。取整数可逆矩阵
\[
V=\begin{pmatrix}-1&3\\1&-2\end{pmatrix},
\qquad \det V=-1,
\]
便有
\[
MV=\begin{pmatrix}1&0\end{pmatrix}.
\]
零列是在两个原关系经过可逆重组后产生的；原关系中却没有一条可以直接删去，因为两条关系合起来定义平凡群，只保留 \(2x=0\) 或只保留 \(3x=0\) 时分别得到 \(C_2\) 或 \(C_3\)。

\paragraph*{求同构类型：读取 Smith 对角元}
例如群由 \(x,y\) 生成，关系为 \(4x+2y=0\)、\(6x+8y=0\)，矩阵为
\[
M=\begin{pmatrix}4&6\\2&8\end{pmatrix}.
\]
一次完整化简为
\[
\begin{pmatrix}4&6\\2&8\end{pmatrix}
\longrightarrow
\begin{pmatrix}2&8\\4&6\end{pmatrix}
\longrightarrow
\begin{pmatrix}2&8\\0&-10\end{pmatrix}
\longrightarrow
\begin{pmatrix}2&0\\0&-10\end{pmatrix}
\longrightarrow
\begin{pmatrix}2&0\\0&10\end{pmatrix}.
\]
四步分别为交换两行、第二行减第一行的两倍、第二列减第一列的四倍、第二行变号。因此 \(A\cong C_2\oplus C_{10}\)。

\paragraph*{追踪原元素：保存行变换矩阵}
若还需要指出商群元素怎样对应，不能只记住对角数字。本例行操作的总矩阵为
\[
U=\begin{pmatrix}0&1\\-1&2\end{pmatrix},\qquad
U^{-1}=\begin{pmatrix}2&-1\\1&0\end{pmatrix}.
\]
原坐标向新坐标转换使用 \(x\mapsto Ux\)；若要把新坐标轴对应的生成元写回原坐标，则使用 \(U^{-1}e_i\)。先在整数坐标上定义同态
\[
\psi:\mathbb Z^2\longrightarrow C_2\oplus C_{10},\qquad
\psi(u,v)=\Bigl([v]_2,[-u+2v]_{10}\Bigr).
\]
前述列操作的矩阵及其整数逆矩阵为
\[
V=\begin{pmatrix}1&-4\\0&1\end{pmatrix},\qquad
V^{-1}=\begin{pmatrix}1&4\\0&1\end{pmatrix}.
\]
实际相乘得到 \(UMV=\operatorname{diag}(2,10)\)。因为 \(V\mathbb Z^2=\mathbb Z^2\)，有 \(U(M\mathbb Z^2)=2\mathbb Z\oplus10\mathbb Z\)。因此
\[
\ker\psi=U^{-1}(2\mathbb Z\oplus10\mathbb Z)=M\mathbb Z^2.
\]
\(U\) 在整数上可逆，随后分别取模的映射满射，所以 \(\psi\) 满射。

\ProseParagraphBreak
由\cref{b1:thm:first-isomorphism}，\(\psi\) 诱导同构 \(\mathbb Z^2/M\mathbb Z^2\rightarrow C_2\oplus C_{10}\)，把陪集 \((u,v)+M\mathbb Z^2\) 送到 \(\psi(u,v)\)。原坐标中的 \(U^{-1}e_1=(2,1)\) 对应 \(C_2\) 的标准生成元，\(U^{-1}e_2=(-1,0)\) 对应 \(C_{10}\) 的标准生成元；它们在商群中的阶分别为二与十。

\ProseParagraphBreak
若生成元为 \(x,y,z\)，关系只有 \(2x=0\)、\(x+3y=0\)，则先用第二条消去 \(x=-3y\)，第一条成为 \(6y=0\)，而 \(z\) 没有关系。因此群为 \(C_6\oplus\mathbb Z\)。没有非零对角关系约束的位置贡献自由生成元，而不是有限阶循环因子。

\subsection{整数格的商群与有限指数}
\label{b1:subsec:0904:02}

\begin{proposition}[子格指数与商群结构]\label{b1:prop:lattice-index}
设 \(m,n\) 为正整数，\(M\) 为 \(m\times n\) 整数矩阵，\(L=M\mathbb Z^n\le\mathbb Z^m\)。若 \(M\) 的 Smith 标准形有 \(t\) 个非零对角元 \(d_1,\ldots,d_t\)，则
\[
\mathbb Z^m/L\cong\mathbb Z^{m-t}\oplus
\bigoplus_{i=1}^t C_{d_i}.
\]
特别地，\(L\) 的指数有限当且仅当 \(t=m\)，此时指数为 \(\prod_i d_i\)。若 \(M\) 为非奇异的 \(m\times m\) 矩阵，则指数为 \(|\det M|\)。
\end{proposition}

\begin{proof}
由\cref{b1:lem:smith-integer}，存在整数可逆矩阵 \(U,V\) 使 \(UMV=D\)，其中 \(D\) 是所述对角阵。由于 \(V\mathbb Z^n=\mathbb Z^n\)，有 \(U(L)=D\mathbb Z^n=\langle d_1e_1,\ldots,d_te_t\rangle\)。因此 \(x+L\mapsto Ux+U(L)\) 给出商群同构。再把前 \(t\) 个坐标分别模 \(d_i\)，其余坐标保留，所得满同态的核恰为 \(U(L)\)；\cref{b1:thm:first-isomorphism}遂给出所述商群公式。

指数就是商群的元素数。当 \(m-t>0\) 时有无限循环直和因子，故指数无限；当 \(m=t\) 时指数为有限循环因子阶的乘积。若 \(M\) 是非奇异方阵，则 \(t=m\)，而 \(\det U,\det V\in\{1,-1\}\)，所以从 \(UMV=D\) 得 \(|\det M|=\det D=\prod_i d_i\)。
\end{proof}

例如上一小节的二维关系子格指数为二十，等于 \(|4\cdot8-6\cdot2|\)。但行列式只给出总阶，并不决定商群：对角矩阵 \(\operatorname{diag}(1,20)\) 与 \(\operatorname{diag}(2,10)\) 的行列式同为二十，商群分别为 \(C_{20}\) 与 \(C_2\oplus C_{10}\)，两者中元素的最大阶不同。

\begin{proposition}[行列式因子]\label{b1:prop:determinantal-divisors}
设 \(m,n\) 为正整数，\(M\) 为 \(m\times n\) 整数矩阵。对 \(1\le k\le\min(m,n)\)，令 \(\Delta_k(M)\) 为 \(M\) 的全部 \(k\) 阶子式的非负最大公因数，并将它称为第 \(k\) 个\term[xinglieshi-yinzi]{行列式因子}[determinantal divisor]；全部相应子式都为零时约定其最大公因数为零，并令 \(\Delta_0(M)=1\)。令
\[
t=\operatorname{rank}_{\mathbb Q}M
\]
为 \(M\) 在有理数域上的矩阵秩。这个数也等于像子群 \(M\mathbb Z^n\) 的自由秩，以及 Smith 标准形中非零对角元的个数。将这些非零对角元依次记为 \(d_1,\ldots,d_t\)，则
\[
\begin{aligned}
\Delta_k(M)&=d_1\cdots d_k &&(1\le k\le t),\\
\Delta_k(M)&=0 &&(t<k\le\min(m,n)).
\end{aligned}
\]
因此比值公式为
\[
d_k=\frac{\Delta_k(M)}{\Delta_{k-1}(M)}\qquad(1\le k\le t).
\]
\end{proposition}

\symidx{\(\Delta_k(M)\)：矩阵的行列式因子}

\begin{proof}
整数可逆矩阵在有理数域上也可逆，所以 Smith 化简 \(UMV=D\) 保持 \(\operatorname{rank}_{\mathbb Q}\)，而 \(D\) 的有理秩就是其非零对角元个数。另一方面，\(U\) 把像子群 \(M\mathbb Z^n\) 同构地送到由 \(d_1e_1,\ldots,d_te_t\) 生成的自由交换群，所以该像子群的自由秩也等于 \(t\)。

整数初等行列操作保持各阶子式的最大公因数。以某行加另一行的整数倍为例，每个受影响子式由行列式的线性性写成原子式与另一原子式的整数线性组合；如果所加行已包含在子式中，额外项有重复行而为零。逆操作给出反向整除，所以最大公因数不变。换行、变号及列操作同样成立。

对 Smith 对角阵，当 \(1\le k\le t\) 时，非零 \(k\) 阶子式是所选 \(k\) 个对角元的乘积。若指标为 \(i_1<\cdots<i_k\)，则 \(i_j\ge j\)，整除链使 \(d_1\cdots d_k\) 整除该乘积；而前 \(k\) 个对角元的子式恰好等于此乘积。因此它就是全部子式的最大公因数。各项为正，结合 \(\Delta_0(M)=1\)，便得到所写的比值公式。

若 \(k>t\)，对角阵的任意 \(k\) 阶子式都为零，故 \(\Delta_k(M)=0\)。特别地，零矩阵的所有正阶行列式因子均为零，不适用上述比值公式。
\end{proof}

\begin{corollary}[满行秩矩形矩阵的子格指数]\label{b1:cor:full-row-rank-index}
设 \(m,n\) 为正整数，\(M\) 为满行秩的 \(m\times n\) 整数矩阵，即 \(\operatorname{rank}_{\mathbb Q}M=m\)。则 \(m\le n\)，并且
\[
[\mathbb Z^m:M\mathbb Z^n]=\Delta_m(M).
\]
即子格 \(M\mathbb Z^n\) 在 \(\mathbb Z^m\) 中的指数等于全部 \(m\) 阶子式绝对值的最大公因数。若 \(m=n\)，这就是 \(|\det M|\)。
\end{corollary}

\begin{proof}
满行秩意味着非零 Smith 对角元的个数 \(t=m\)。又因 \(t\leq\min(m,n)\)，必有 \(m\leq n\)。由\cref{b1:prop:lattice-index}，指数为 \(d_1\cdots d_m\)；再由\cref{b1:prop:determinantal-divisors}，此乘积等于 \(\Delta_m(M)\)。方阵只有一个 \(m\) 阶子式，即其行列式。
\end{proof}

这里必须是满行秩，而不能笼统地只说矩形矩阵“满秩”。若 \(n<m\) 且 \(M\) 满列秩，则 \(t=n<m\)，商群仍有自由秩 \(m-n\)，因而子格指数无限；此时全部最大非零子式是 \(n\) 阶子式，它们的最大公因数并不表示 \(M\mathbb Z^n\) 在 \(\mathbb Z^m\) 中的有限指数。

\ProseParagraphBreak
前面用于说明零列的矩阵 \(M=\begin{pmatrix}2&3\end{pmatrix}\) 有满行秩，而
\[
M\mathbb Z^2=2\mathbb Z+3\mathbb Z=\mathbb Z,
\]
因为 \(3-2=1\)。其指数为 \(1=\gcd(2,3)\)，不能任选一个非零最大阶子式 \(2\) 或 \(3\) 作为指数。

\ProseParagraphBreak
行列式因子也适合核验小矩阵的 Smith 结果。本节的二维算例中，全部元素的最大公因数为二，唯一二阶子式为二十，因此 \(d_1=2,d_2=10\)。大矩阵的子式很多，可以先做整数行列化简，再抽查少量子式以发现计算失误。不过，若所抽取子式的最大公因数为 \(g\)，只能直接得到 \(\Delta_k(M)\mid g\)，不能据此断言相等。例如 \(\operatorname{diag}(2,3)\) 的全部一阶子式最大公因数为 \(1\)，只检查条目 \(2\) 就不能恢复它。

\ProseParagraphBreak
若要独立确认正的候选值 \(\delta=\Delta_k(M)\)，须证明全部 \(k\) 阶子式均被 \(\delta\) 整除，并找到最大公因数恰为 \(\delta\) 的一组子式。这两个方向合起来才给出等号。另一种充分验证是完整保存 \(U,V\)，核验 \(UMV=D\) 及 \(U,V\) 的整数逆矩阵，并检查 \(D\) 的非零对角元为正且形成整除链；这样便确认了 Smith 结果，无须枚举全部子式。

\subsection{模论解释与后续推广预告}
\label{b1:subsec:0904:03}

本章始终只用整数倍、子群与商群。下一卷将把这种结构重新解释为模：整数 \(n\) 对交换群元素 \(x\) 的作用就是 \(nx\)。这个预告不作为当前证明的输入；模的定义与公理将在第二卷第二章“模、子模与商模”正式建立。

\ProseParagraphBreak
从这一后续观点看，\(\mathbb Z^m\) 是有限秩自由对象，关系矩阵记录了从一个自由对象到另一个自由对象的映射，商群记录生成元在关系约束下留下的信息。整数的任意理想由一个元素生成，这一性质将被抽象为主理想整环条件，从而把本章的循环分解推广到更一般的系数系统。在本章，支持消元的具体工具是整数除法以及 Bézout 等式；推广时必须重新核对系数系统是否提供相应工具。

\ProseParagraphBreak
例如对有理数系数多项式也可以作带余除法，第二卷第四章“主理想整环上的模与标准形”会用类似方法研究线性变换的有理标准形。届时矩阵中的“整数关系”变为“多项式关系”，循环因子表达一个向量及其在反复作用下的生成规律。这个类比解释了有限生成交换群为什么安排在群论之后：它既完成一项独立分类，也提供了从群的生成关系过渡到模的生成关系的具体经验。

\ProseParagraphBreak
实际使用本章结果时，可按三步整理资料。第一步，确认群交换且有限生成，选择有限生成元，并取得一张列向量生成整个关系核的有限整数矩阵。第二步，明确行列约定，在整数上作可逆变换，同时累计 \(U,V\)，直至得到 \(D=UMV\)。第三步，核验 \(U,V\) 都有整数逆矩阵，且 \(D\) 的非零对角元为正并形成整除链，再从 \(D\) 读取自由秩及非平凡不变因子。

\ProseParagraphBreak
有限秩子群定理只保证完整关系核存在有限基，并不自动给出可供计算的生成组。若矩阵只记录部分已知关系，只能得到相应展示群到原群的自然满同态，不能把展示群的分类结果直接当作原群的分类。若需要追踪特定元素或子群，应使用保存的坐标变换；若只需判断两个有限生成交换群是否同构，比较自由秩与归一化后的不变因子即可。


\begin{chaptersummary}
有限生成阿贝尔群的关系核是有限秩自由阿贝尔群，因此可用有限整数矩阵表示。Smith 化简只作可逆整数行列操作，得到自由部分与按整除次序排列的循环因子。自由秩由无挠商群确定，每个素数幂因子的长度由 \(p^{j-1}A_p/p^jA_p\) 的阶确定。矩阵行列式给出满秩商格的指数，各阶子式的最大公因数进一步恢复全部不变因子；只有总阶通常不足以区分群。
\end{chaptersummary}
\BookExercises

\begin{exercise}\label{b1:ex:09-congruence-lattice}
求子群
\[
H=\{\,(x,y,z)\in\mathbb Z^3\mid
x+2y+3z\equiv0\pmod6\,\}
\]
的一组基，以及 \(\mathbb Z^3/H\) 的结构。
\end{exercise}
\begin{solution}
每个元素唯一写为
\[
k(6,0,0)+y(-2,1,0)+z(-3,0,1),
\qquad k=(x+2y+3z)/6\in\mathbb Z.
\]
因此所列三个向量为一组基。同态
\(\mathbb Z^3\to C_6\)、\((x,y,z)\mapsto\overline{x+2y+3z}\)
满射且核为 \(H\)，故商群为 \(C_6\)，指数为六。
\end{solution}

\begin{exercise}\label{b1:ex:09-explicit-quotient}
求 \(A=\mathbb Z^2/\langle(6,0),(4,10)\rangle\) 的不变因子，并在原坐标中选出对应两个循环直和因子的生成元。
\end{exercise}
\begin{solution}
关系矩阵为 \(\begin{pmatrix}6&4\\0&10\end{pmatrix}\)。一阶子式最大公因数为二，行列式为六十，所以不变因子为 \(2,30\)。令 \(u=(3,0)+L\)、\(v=(0,1)+L\)，其中 \(L\) 为关系格。\(u\) 的阶为二；若 \(kv=0\)，则
\[
(0,k)=a(6,0)+b(4,10),
\]
所以 \(k=10b\)、\(6a+4b=0\)，迫使 \(3\mid b\)，最小正 \(k\) 为三十。若 \(u\) 属于 \(\langle v\rangle\)，某个关系向量的第一坐标应为三，但所有关系向量的第一坐标均为偶数，不可能。故两个循环子群交为零，乘积有六十个元素，等于全群，给出 \(A=\langle u\rangle\oplus\langle v\rangle\)。
\end{solution}

\begin{exercise}\label{b1:ex:09-order-360}
列出所有阶为 \(360\) 的阿贝尔群，写成不变因子形式，并说明没有遗漏或重复。
\end{exercise}
\begin{solution}
\(360=2^3\cdot3^2\cdot5\)。二主分量对应三的分拆 \(3,\ 2+1,\ 1+1+1\)，三主分量对应二的分拆 \(2,\ 1+1\)，五主分量只能为 \(C_5\)。共六种组合，按素数指数左端补零并逐列相乘，得到
\[
\begin{gathered}
C_{360},\qquad C_3\oplus C_{120},\qquad
C_2\oplus C_{180},\\
C_6\oplus C_{60},\qquad
C_2\oplus C_2\oplus C_{90},\qquad
C_2\oplus C_6\oplus C_{30}.
\end{gathered}
\]
每条整除链都满足不变因子条件，阶均为三百六十。主分解及素数幂循环因子的唯一性保证没有遗漏，互不相同的不变因子链保证没有重复。
\end{solution}

\begin{exercise}\label{b1:ex:09-free-quotient-splitting}
设 \(B\le A\)，其中 \(A\) 阿贝尔，且 \(A/B\) 为自由阿贝尔群。证明存在子群 \(C\) 使 \(A=B\oplus C\)。举例说明 \(C\) 未必唯一。
\end{exercise}
\begin{solution}
对 \(A/B\) 的每个基元素选一个提升 \(c_i\in A\)。自由阿贝尔群的泛性质给出同态 \(s:A/B\to A\)，把基元素送到对应提升；商映射 \(\pi\) 满足 \(\pi s=1\)。令 \(C=s(A/B)\)。任意 \(a\in A\) 可写成 \((a-s\pi(a))+s\pi(a)\)，首项属于 \(B\)，故 \(A=B+C\)。若 \(s(x)\in B\)，施以 \(\pi\) 得 \(x=0\)，故交为零。对 \(A=\mathbb Z\oplus C_2\)、\(B=0\oplus C_2\)，子群 \(\langle(1,\bar0)\rangle\) 与 \(\langle(1,\bar1)\rangle\) 都是这样的补群，且彼此不同。
\end{solution}

\begin{exercise}\label{b1:ex:09-cyclic-hom}
把同态集合按逐点相加组成阿贝尔群，证明
\[
\operatorname{Hom}(C_m,C_n)\cong C_{\gcd(m,n)}.
\]
进而计算从 \(C_{12}\oplus C_{18}\) 到 \(C_6\) 的同态数与满同态数。
\end{exercise}
\begin{solution}
同态由 \(\bar1\) 的像 \(x\in C_n\) 决定，条件是 \(mx=0\)。令 \(d=\gcd(m,n)\)，满足条件的元素恰为 \(\overline{(n/d)t}\)，其中 \(t\) 模 \(d\) 取值，在加法下构成 \(C_d\)。本题两生成元在 \(C_6\) 中都可以任取像，所以共有三十六个同态。若像由 \(\bar a,\bar b\) 生成，满射当且仅当 \(\gcd(a,b,6)=1\)。非满情况是两者都为二的倍数或都为三的倍数：分别有九种与四种，两者都满足有一种。故满同态数为 \(36-9-4+1=24\)。
\end{solution}

\begin{exercise}\label{b1:ex:09-minimal-generators}
设
\[
A\cong\mathbb Z^r\oplus C_{d_1}\oplus\cdots\oplus C_{d_s},
\qquad 1<d_1\mid\cdots\mid d_s.
\]
证明 \(A\) 的最少生成元数为 \(r+s\)。
\end{exercise}
\begin{solution}
各因子的标准生成元给出 \(r+s\) 个生成元。若 \(s>0\)，选素数 \(p\mid d_1\)，则 \(p\) 整除全部 \(d_i\)，所以 \(A/pA\cong C_p^{\,r+s}\)。任何 \(k\) 个生成元的像至多产生 \(p^k\) 个元素，因为每个整数系数可模 \(p\) 约减；而该商群有 \(p^{r+s}\) 个元素，故 \(k\ge r+s\)。若 \(s=0\)，任选素数 \(p\)，同样用 \(A/pA\cong C_p^r\) 得到下界 \(r\)。单位群对应 \(r=s=0\)，无需生成元。
\end{solution}

\begin{exercise}\label{b1:ex:09-torsion-layers}
设有限阿贝尔二群 \(A\) 满足
\[
|A/2A|=8,\quad |2A/4A|=4,\quad
|4A/8A|=2,\quad 8A=0.
\]
确定 \(A\)，并求其中阶恰为四的元素个数。
\end{exercise}
\begin{solution}
指数至少为一、二、三的循环因子分别有三、二、一个，因此指数恰为一、二、三的各有一个，故
\(A\cong C_2\oplus C_4\oplus C_8\)。被四消去的元素数为 \(2\cdot4\cdot4=32\)，被二消去的元素数为 \(2\cdot2\cdot2=8\)。阶恰为四的元素就是两者之差，共二十四个。
\end{solution}

\begin{exercise}\label{b1:ex:09-fg-endomorphism}
证明有限生成阿贝尔群的每个满自同态都是同构，并举出一个单自同态不是满射的例子。有限生成条件能否从第一个结论中删除？
\end{exercise}
\begin{solution}
设 \(f:A\to A\) 满射，\(T=A_{\mathrm{tor}}\)。它诱导 \(\bar f:A/T\to A/T\) 的满自同态。写 \(A/T\cong\mathbb Z^r\)，其整数矩阵 \(M\) 有整数右逆：为每个标准基选原像并列为矩阵 \(N\)，得 \(MN=I\)。行列式给出 \(\det M=\pm1\)，所以 \(\bar f\) 为同构。于是若 \(f(a)\in T\)，则 \(a\in T\)。满射性因而限制为 \(f|_T:T\to T\) 的满射；\(T\) 有限，所以这个限制为同构。若 \(f(a)=0\)，先得 \(a\in T\)，再得 \(a=0\)，故 \(f\) 同构。

单而不满的例子为 \(\mathbb Z\to\mathbb Z\)、\(n\mapsto2n\)。有限生成条件不能删除：在 \(\bigoplus_{i\ge1}C_2\) 上，把第一个标准生成元送零，把第 \(i+1\) 个送到第 \(i\) 个，便是有非平凡核的满自同态。
\end{solution}
